% ── §2.4: Ψηφιακός Δίδυμος Χρήστη (Digital Twin) ────────────

\section{Ψηφιακός Δίδυμος Χρήστη (Digital Twin)}
\label{sec:digital_twin}

Η έννοια του Ψηφιακού Διδύμου (Digital Twin, DT) ορίζεται ως μια
δυναμική ψηφιακή αναπαράσταση ενός φυσικού οντότητας ή συστήματος,
η οποία ενημερώνεται συνεχώς από ροές δεδομένων πραγματικού χρόνου
\cite{Grieves2014}. Αρχικά αναπτύχθηκε στον χώρο της βιομηχανικής
παραγωγής, αλλά η εφαρμογή του στη ψηφιακή υγεία και τα AAL
συστήματα αναδεικνύεται ως πεδίο υψηλής ερευνητικής δραστηριότητας
\cite{Tao2019}.

\subsection{Στρώματα Αρχιτεκτονικής}
\label{subsec:dt_layers}

Ένας Ψηφιακός Δίδυμος χρήστη σε AAL περιβάλλον οργανώνεται σε τρία
εννοιολογικά στρώματα \cite{LiuDT2021}:

\begin{enumerate}
    \item \textbf{Στρώμα Δεδομένων (Data Layer):} Συλλέγει ακατέργαστα
    δεδομένα αισθητήρων από το smartphone (επιταχυνσιόμετρο, γυροσκόπιο)
    και άλλες πηγές (wearables, έξυπνο σπίτι). Αποτελεί την \textit{«αίσθηση»}
    του διδύμου — το τι συμβαίνει στον φυσικό κόσμο.

    \item \textbf{Στρώμα Μοντέλου (Model Layer):} Εκτελεί επεξεργασία
    και ταξινόμηση των δεδομένων μέσω μοντέλων μηχανικής μάθησης.
    Στην παρούσα εργασία, το στρώμα αυτό υλοποιείται από το 1D-CNN
    μοντέλο HAR που εκπαιδεύεται συνεργατικά μέσω FL.

    \item \textbf{Στρώμα Υπηρεσιών (Service Layer):} Παρέχει ανθρώπινα
    αναγνώσιμες πληροφορίες — ειδοποιήσεις, τάσεις δραστηριότητας,
    αναφορές σε φροντιστές — βάσει των εξόδων του στρώματος μοντέλου.
\end{enumerate}

\subsection{Ψηφιακός Δίδυμος στο FL Σύστημα}
\label{subsec:dt_in_fl}

Στο πλαίσιο της παρούσας εργασίας, κάθε client του FL συστήματος
αντιστοιχεί σε έναν χρήστη του AAL περιβάλλοντος. Το τοπικό μοντέλο
HAR κάθε client λειτουργεί ως \textit{εξατομικευμένος Ψηφιακός Δίδυμος}:
αποτυπώνει τα ιδιαίτερα κινηματικά μοτίβα του συγκεκριμένου χρήστη
χωρίς να μεταδίδει ευαίσθητα δεδομένα εκτός της συσκευής.

Η έννοια του Ψηφιακού Διδύμου επιβεβαιώνει τη σχεδιαστική επιλογή
να διατηρούμε τοπικά δεδομένα στις συσκευές: ο DT κάθε χρήστη
\textit{ζει} στη συσκευή του, ενώ η γνώση (μοντέλο) μεταδίδεται στον
server μέσω FL. Αυτή η αρχιτεκτονική εκμεταλλεύεται τα πλεονεκτήματα
τόσο του FL (προστασία ιδιωτικότητας) όσο και του DT (εξατομίκευση
και συνεχής ενημέρωση). Η συνολική αρχιτεκτονική DT+FL αποτυπώνεται
στο Σχήμα~\ref{fig:dt_fl_arch}.

\begin{figure}[H]
\centering
\begin{tikzpicture}[
    node distance=1cm and 1.8cm,
    sensor/.style={rectangle, rounded corners=3pt, draw=gray!70, thick,
                   fill=gray!10, minimum width=2cm, minimum height=0.7cm,
                   font=\scriptsize, align=center},
    dtbox/.style={rectangle, rounded corners=5pt, draw=blue!60, thick,
                  fill=blue!8, minimum width=3.2cm, minimum height=0.8cm,
                  font=\small, align=center},
    srvbox/.style={rectangle, rounded corners=5pt, draw=green!60!black, thick,
                   fill=green!10, minimum width=3.5cm, minimum height=0.9cm,
                   font=\small, align=center},
    arrow/.style={-{Stealth[length=6pt]}, thick},
    dbarrow/.style={<->, thick, dashed, gray!70}
]

% Physical world (left column)
\node[sensor] (phone1) {Smartphone 1\\[2pt](sensors)};
\node[sensor, below=0.5cm of phone1] (phone2) {Smartphone 2\\[2pt](sensors)};
\node[font=\scriptsize, below=0.3cm of phone2] (dots) {$\vdots$};
\node[sensor, below=0.3cm of dots] (phoneN) {Smartphone $N$\\[2pt](sensors)};

% DT nodes (middle column)
\node[dtbox, right=2.2cm of phone1] (dt1) {DT User 1\\[2pt](local 1D-CNN)};
\node[dtbox, right=2.2cm of phone2] (dt2) {DT User 2\\[2pt](local 1D-CNN)};
\node[font=\small, right=2.2cm of dots] (dots2) {$\vdots$};
\node[dtbox, right=2.2cm of phoneN] (dtN) {DT User $N$\\[2pt](local 1D-CNN)};

% Server (right)
\node[srvbox, right=2.5cm of dt2] (srv) {FL Server\\[2pt]FedAvg + Incentives};

% Sensor → DT (raw data stays local)
\draw[dbarrow] (phone1) -- node[above, font=\tiny]{raw data} (dt1);
\draw[dbarrow] (phone2) -- node[above, font=\tiny]{raw data} (dt2);
\draw[dbarrow] (phoneN) -- node[above, font=\tiny]{raw data} (dtN);

% DT → Server (only model weights)
\draw[arrow, blue!60] (dt1.east) -- node[above, font=\tiny]{$\Delta\mathbf{w}_1$} (srv.west|-dt1);
\draw[arrow, blue!60] (dt2.east) -- node[above, font=\tiny]{$\Delta\mathbf{w}_2$} (srv.west|-dt2);
\draw[arrow, blue!60] (dtN.east) -- node[above, font=\tiny]{$\Delta\mathbf{w}_N$} (srv.west|-dtN);

% Server → DT (global model)
\draw[arrow, green!60!black] (srv.west|-dt1) -- +(-0.4,0) -- (dt1.east);

% Labels
\node[font=\tiny, text=gray, below=0.15cm of dt1] {Τοπική εκπαίδευση};
\node[font=\tiny, text=red!60!black, above=0.15cm of phone1] {Δεδομένα: τοπικά};
\node[font=\tiny, text=green!60!black, above right=0.1cm and -0.3cm of srv]
    {Global model};

\end{tikzpicture}
\caption{Αρχιτεκτονική DT+FL: κάθε χρήστης AAL αντιστοιχεί σε έναν Ψηφιακό Δίδυμο που αποθηκεύει τοπικά τα ακατέργαστα δεδομένα αισθητήρων και στέλνει στον FL server μόνο ενημερώσεις βαρών μοντέλου ($\Delta\mathbf{w}_i$). Ο server συγκεντρώνει, αξιολογεί τη συνεισφορά και κατανέμει ανταμοιβές~\cite{LiuDT2021, Tao2019}.}
\label{fig:dt_fl_arch}
\end{figure}
